\documentclass[12pt, twoside]{article}
%\documentclass[12pt, twoside]{article}
\usepackage[letterpaper, margin=1in, headsep=0.2in]{geometry}
\setlength{\headheight}{0.6in}
%\usepackage[english]{babel}
\usepackage[utf8]{inputenc}
\usepackage{microtype}
\usepackage{amsmath}
\usepackage{amssymb}
%\usepackage{amsfonts}
\usepackage[nomessages]{fp} %\FPeval{\var-name}{2*sin(pi/6)}
\usepackage{siunitx} %units in math. eg 20\milli\meter
\usepackage{yhmath} % for arcs, overparenth command
\usepackage{tikz} %graphics
\usetikzlibrary{quotes, angles, arrows, arrows.meta}
\usepackage{pgfplots}
\usepgfplotslibrary{fillbetween}
\pgfplotsset{compat=1.16}
\usepackage{graphicx} %consider setting \graphicspath{{images/}}
\usepackage{parskip} %no paragraph indent
\usepackage{enumitem}
\usepackage{multicol}
\usepackage{venndiagram}

\usepackage{fancyhdr}
\pagestyle{fancy}
\fancyhf{}
\renewcommand{\headrulewidth}{0pt} % disable the underline of the header
\raggedbottom
\hfuzz=2mm %suppresses overfull box warnings

\usepackage{hyperref}

\fancyhead[LO]{La Scuola d'Italia / Dr.\ Huson / IB Math 11 \\* 7 October 2026}
\fancyhead[RO]{\fbox{\begin{minipage}{6cm}\footnotesize Nickname:\\[0.5cm]\end{minipage}}}
\fancyhead[LE]{\fbox{\begin{minipage}{6cm}\footnotesize Nickname:\\[0.5cm]\end{minipage}}}
\fancyfoot[C]{\thepage}

\begin{document}
\subsubsection*{1.10 Quiz: Arithmetic and Geometric Sequences}

\emph{A calculator is allowed. Show your working: write the formula you
use, then substitute.}

\begin{enumerate}[itemsep=0.15cm]

%--- Problem 1: Arithmetic sequence; d, u_n, u_30, show 303 is a term [8 marks] ---
%--- Source: Original (freshly authored, AA SL 1.2); parallel of 1-9PQ_Sequences P1(a)-(d) ---
\item[1.] [Maximum mark: 8]

Consider the arithmetic sequence
$$15,\; 23,\; 31,\; 39,\; \ldots$$

\begin{enumerate}[itemsep=0.1cm]
  \item Write down the value of the common difference $d$. \hfill [1]
  \item Find an expression for $u_{n}$, the $n^{\text{th}}$ term. \hfill [2]
  \item Find $u_{30}$. \hfill [2]
  \item Show that $303$ is a term of the sequence, and find which term
  it is. \hfill [3]
\end{enumerate}

\begin{tikzpicture}
  \draw (-0.6,0) rectangle (14.6,14.6);
  \foreach \y in {1,2,...,14} \draw[dotted] (0,\y)--(14.1,\y);
\end{tikzpicture}

\newpage

%--- Problem 2: Two given terms -> u_1 and d; first negative term; terms in k [9 marks] ---
%--- Source: Original (freshly authored, AA SL 1.2); parallel of 1-9PQ_Sequences P2(a), (c), (d) ---
\item[2.] [Maximum mark: 9]

In an arithmetic sequence, $u_{5} = 38$ and $u_{14} = 11$.

\begin{enumerate}[itemsep=0.1cm]
  \item Find the common difference $d$ and the first term $u_{1}$.
  \hfill [3]
  \item Find the first term of the sequence that is negative. \hfill [3]
\end{enumerate}

The first three terms of a different arithmetic sequence are
$$2k + 1, \qquad 4k - 3, \qquad 5k + 1.$$

\begin{enumerate}[itemsep=0.1cm]
  \item[(c)] Find the value of $k$, and write down the first three terms.
  \hfill [3]
\end{enumerate}

\begin{tikzpicture}
  \draw (-0.6,0) rectangle (14.6,16.3);
  \foreach \y in {1,2,...,16} \draw[dotted] (0,\y)--(14.1,\y);
\end{tikzpicture}

\newpage

%--- Problem 3: Arithmetic, geometric or neither; geometric r and u_12; challenge: two given terms [10 marks] ---
%--- Source: Original (freshly authored, AA SL 1.2, 1.3); parallel of 1-9PQ_Sequences P3(a), (b), (d), (f) ---
\item[3.] [Maximum mark: 10]

\begin{enumerate}[itemsep=0.1cm]
  \item State whether each sequence is arithmetic, geometric or neither.
  Write down $d$ or $r$ where there is one. \hfill [4]

  \hspace*{1em} (i) $5,\; 10,\; 20,\; 40$ \hfill (ii) $7,\; 4,\; 1,\; -2$
  \hfill (iii) $2,\; 3,\; 5,\; 8$ \hfill (iv) $162,\; 54,\; 18,\; 6$
  \hspace*{1em}
\end{enumerate}

Consider the geometric sequence
$$16,\; 20,\; 25,\; \ldots$$

\begin{enumerate}[itemsep=0.1cm]
  \item[(b)] Find the common ratio $r$. \hfill [1]
  \item[(c)] Find $u_{12}$. \hfill [2]
  \item[(d)] Challenge: A different geometric sequence has $u_{3} = 45$
  and $u_{6} = 1215$. Find the common ratio and the first term. \hfill [3]
\end{enumerate}

\begin{tikzpicture}
  \draw (-0.6,0) rectangle (14.6,14.4);
  \foreach \y in {1,2,...,14} \draw[dotted] (0,\y)--(14.1,\y);
\end{tikzpicture}

\newpage

%--- Problem 4: Arithmetic vs geometric salary; expressions, year 3, first year B > A [7 marks] ---
%--- Source: Original (freshly authored, AA SL 1.2, 1.3); parallel of 1-9PQ_Sequences P4(a)-(c), (a) and (b) swapped ---
\item[4.] [Maximum mark: 7]

Marco is offered two jobs.

\begin{center}
\begin{tabular}{rl}
Job A: & \$$38\,000$ in the first year, rising by \$$2000$ each year. \\[2pt]
Job B: & \$$36\,000$ in the first year, rising by $6\%$ each year. \\
\end{tabular}
\end{center}

\begin{enumerate}[itemsep=0.1cm]
  \item Write down an expression for the salary in year $n$ for each job.
  \hfill [2]
  \item Find Marco's salary in year $3$ for each job. \hfill [2]
  \item Find the first year in which the Job B salary is greater than the
  Job A salary. \hfill [3]
\end{enumerate}

\begin{tikzpicture}
  \draw (-0.6,0) rectangle (14.6,16.5);
  \foreach \y in {1,2,...,16} \draw[dotted] (0,\y)--(14.1,\y);
\end{tikzpicture}

\end{enumerate}
\end{document}
